% Fixture: a single-column article with plain equation numbers, numeric citations,
% amsthm statements and proofs, and no PDF outline. Compile with: tectonic article-numeric.tex
\documentclass[11pt]{article}
\usepackage{amsmath,amssymb,amsthm}
\usepackage[margin=1.2in]{geometry}
\newtheorem{theorem}{Theorem}
\newtheorem{lemma}[theorem]{Lemma}
\newtheorem{corollary}[theorem]{Corollary}
\theoremstyle{definition}
\newtheorem{definition}{Definition}
\begin{document}
\title{Shrinkage and Sparsity in Linear Models}
\author{A. Fixture}
\date{}
\maketitle

\section{Introduction}
We study the linear model in which the response depends linearly on $p$ predictors.
Least squares is the classical estimator and was analysed in detail by many authors [1, 3].
Throughout, let $X$ denote the $N \times p$ design matrix and let $y$ denote the response vector.
The residual sum of squares is
\begin{equation}
\mathrm{RSS}(\beta) = \sum_{i=1}^{N} (y_i - x_i^T \beta)^2 .
\label{eq:rss}
\end{equation}
Minimizing \eqref{eq:rss} over $\beta$ gives the normal equations, which are discussed in Section~2.

\section{Ridge regression}
The ridge estimator adds a penalty to the criterion in \eqref{eq:rss}. Formally,
\begin{equation}
\hat\beta^{\mathrm{ridge}} = \operatorname*{argmin}_{\beta} \Big\{ \mathrm{RSS}(\beta) + \lambda \sum_{j=1}^{p} \beta_j^2 \Big\},
\label{eq:ridge}
\end{equation}
where $\lambda \ge 0$ is a complexity parameter that controls the amount of shrinkage [2].
\begin{definition}
The \emph{effective degrees of freedom} of the ridge fit is the trace of the hat matrix.
\end{definition}
\begin{lemma}\label{lem:solution}
For every $\lambda > 0$ the minimizer in \eqref{eq:ridge} is unique and equals $(X^TX + \lambda I)^{-1} X^T y$.
\end{lemma}
\begin{proof}
The criterion in Eq.~\eqref{eq:ridge} is strictly convex when $\lambda > 0$, so its minimizer is unique.
Setting the gradient to zero gives the stated formula.
\end{proof}
\begin{theorem}[Shrinkage]\label{thm:shrink}
The ridge coefficients shrink towards zero as $\lambda$ grows. In particular, the norm of the estimator in Lemma~\ref{lem:solution} is a decreasing function of $\lambda$.
\end{theorem}
\begin{proof}
Write the singular value decomposition $X = UDV^T$. By Lemma~\ref{lem:solution},
\begin{equation}
X\hat\beta^{\mathrm{ridge}} = \sum_{j=1}^{p} u_j \frac{d_j^2}{d_j^2 + \lambda} u_j^T y ,
\label{eq:svd}
\end{equation}
and each factor in \eqref{eq:svd} decreases in $\lambda$, see [3, Theorem 2].
\end{proof}
Figure~\ref{fig:path} shows the coefficient profiles. Table~\ref{tab:df} lists degrees of freedom.

\begin{figure}[h]
\centering
\fbox{\rule{0pt}{1.1in}\hspace{3in}}
\caption{Profiles of ridge coefficients as the penalty $\lambda$ varies.}
\label{fig:path}
\end{figure}

\begin{table}[h]
\centering
\caption{Effective degrees of freedom for several penalties.}
\label{tab:df}
\begin{tabular}{cc}
$\lambda$ & df \\ \hline
0 & 8.0 \\
1 & 6.2 \\
10 & 3.1
\end{tabular}
\end{table}

\begin{corollary}
When the predictors are orthonormal, ridge regression rescales the least squares coefficients by $1/(1+\lambda)$, as \eqref{eq:svd} shows.
\end{corollary}

\begin{thebibliography}{9}
\bibitem{hoerl} A. E. Hoerl and R. W. Kennard. Ridge regression: biased estimation for nonorthogonal problems. Technometrics, 12(1):55--67, 1970.
\bibitem{tib} R. Tibshirani. Regression shrinkage and selection via the lasso. Journal of the Royal Statistical Society B, 58(1):267--288, 1996.
\bibitem{htf} T. Hastie, R. Tibshirani, and J. Friedman. The Elements of Statistical Learning. Springer, second edition, 2009.
\end{thebibliography}
\end{document}
